% legends.tex (국문 v10): 그림 설명. 패널마다 사용 연도를 적는다.
\section*{그림 설명}

\textbf{그림 1. ALT 관측, 입력 자료와 그 기간.} (a) ESA CCI 영구동토 확률(1997--2021년 평균)과 세 연구 지역, 2010--2024년 ALT 라벨이 있는 1 km 셀 수. (b) 지역별 관측 연도당 셀-연도 라벨 수(2010--2024년). (c) 자료 기간: 지역별 라벨, 관측 연도 기후 입력에 쓴 ERA5-Land 월자료(2010--2024년), 정적 공변량(연도 없음), 공개 제품 6종의 제공 기간(옅은 띠)과 라벨 셀에서 비교에 쓴 연도(진한 띠). 회색 음영은 분석 기간 2010--2024년이다. (d) 2014년 알래스카 CALM 셀 하나(북위 71.3$^\circ$, 서경 156.6$^\circ$)에서의 기후 입력 정의: ERA5-Land 누적 융해 도일(TDD), 관측일까지 누적한 값(학습 · 채점에 쓰는 관측 연도 입력; 날짜가 없는 CALM 계절 말 값의 관례인 9월 1일), 같은 해의 계절 전체 값(지도에 사용), 2011--2024년 평균(평년 기후 입력).

\textbf{그림 2. 물리 잔차 워크플로와 검증 설계.} (a) 물리 잔차 모델의 구조. 2014년 알래스카 입력을 예로 들었다. 관측 연도 기후 변수와 정적 지형 · 토양 공변량이 재보정 Stefan 모델($\mathrm{ALT}=E\sqrt{\mathrm{TDD}}$) 잔차의 gradient boosting 모델에 들어가고, 관측 ALT는 학습에만 쓴다(파선). 출력은 Stefan 추정값에 잔차의 $\lambda=0.5$배를 더한 값과 적용 가능 영역(2011--2024년 평균 기후로 계산)이다. (b) 비교한 네 구조: 직접 ML, 물리 입력 ML(Stefan · Kudryavtsev 예측을 입력으로), 물리 유사라벨(Stefan 예측으로 학습 자료 증강), 물리 + 잔차 ML. (c) 레나 델타 라벨 셀(2010--2022년)에 표시한 다섯 검증 설계: 무작위, 공간 블록(0.5$^\circ$), 연도 홀드아웃, 블록 × 연도, 지역 홀드아웃(알래스카 · 캐나다로 학습). 설계마다 예시 겹 하나를 표시하였고, 셀-연도 가운데 하나라도 시험 겹에 있으면 시험 셀(검정)로 그렸다.

\textbf{그림 3. 관측 연도 기후 입력, 검증 설계, 잔차 이득의 조건.} 모든 패널은 2010--2024년 셀-연도 라벨과 겹 밖 예측을 쓴다. (a) 평년 기후 입력(2011--2024년 평균, 빈 표식)과 관측 연도 기후 입력(관측일까지 누적, 채운 표식)의 공간 블록 RMSE. 재보정 Stefan 모델, 물리 잔차 모델, 직접 ML. 이 패널에서만 빈 표식과 채운 표식이 기후 입력을 뜻한다. (b) 다섯 검증 설계의 RMSE. 원천 계수 Stefan, 재보정 Stefan, 물리 잔차 모델, 직접 ML. 연도 홀드아웃과 블록 × 연도는 셀-연도 20개 이상인 해를 쓰고, 지역 홀드아웃에서는 두 Stefan 모델이 같으며, 축 밖 값은 숫자로 적었다. (c) 공변량 묶음 여섯 가지에서 물리 잔차 모델과 재보정 Stefan 모델의 공간 블록 RMSE 차이. (d) 여덟 공변량의 지역 안 3분위에서의 같은 차이. 검은 테두리는 95\% 구간이 0을 제외하는 칸, 빗금은 셀-연도 30개 또는 블록 5개 미만인 칸이며, 0.0은 차이가 0.05 cm 미만임을 뜻한다. (a), (c), (d)의 막대와 구간은 0.5$^\circ$ 블록 2,000회 재추출의 95\% 부트스트랩 구간이고, 음수는 물리 잔차 모델이 낫다는 뜻이다.

\textbf{그림 4. 적은 라벨과 새 지역에서의 정확도.} (a) 지역 안 라벨 수 사다리: 2010--2024년 셀-연도에서 뽑은 같은 $n$개 학습 라벨(3회 추출)로 재보정한 Stefan 모델 대비 공간 블록 RMSE 차이. 띠는 물리 잔차 모델의 95\% 블록 부트스트랩 구간이다. (b) 새 지역 전이: 다른 두 지역과 대상 지역 라벨 $n$개로 학습한 모델의 대상 지역 공간 블록 겹 RMSE. (c) 외부 CALM 지역(연구 지역 밖 2010--2024년 관측): 같은 외부 지역에서 더한 라벨 수에 따른 RMSE. Wei 등, Liu 등, ESA CCI v5 값이 모두 있는 7개 지역 \nftrows{}개 셀-연도(2010--2020년)에서 계산하였다. 굵고 옅은 선은 같은 라벨로 재보정한 Stefan 모델로 물리 잔차 모델과 거의 겹친다. 가로선은 제품이며, Wei 등과 Liu 등은 이 지점들로 학습하였고 ESA CCI v5는 ALT 관측으로 학습하지 않았다.

\textbf{그림 5. 같은 셀 · 같은 해의 공개 제품 비교.} (a) Wei 등, Liu 등, ESA CCI v4가 모두 관측 연도 값을 가진 셀-연도(\nyrscc{}년)에서 세 방법(공간 블록 겹 밖 예측)과 제품 5종의 RMSE. 채운 표식은 모든 공통 셀-연도, 빈 표식은 Wei 등의 학습 지점에서 25 km 넘게 떨어진 셀만 쓴 값이며, ESA CCI v5와 Aalto 등은 공통 셀-연도 가운데 그 제품 값이 있는 셀-연도에서 채점하였고, 축 밖 값은 숫자로 적었다. (b) 관측 연도별 RMSE. 해마다 같은 셀-연도에서 물리 잔차 모델, 재보정 Stefan 모델, Wei 등, Liu 등(2020년까지)을 계산하였고, 셀-연도 30개 이상인 해(알래스카 \nyrsywAkko, 레나 델타 \nyrsywLeko, 캐나다 \nyrsywCako)만 표시하였다. 회색 막대는 해마다의 셀-연도 수(로그 축)이다.

\textbf{그림 6. 연도가 명시된 1 km 지도.} (a) 알래스카 물리 잔차 모델의 2011--2024년 평균 ALT. 연도별 지도의 평균이며, 각 연도 지도는 그 해의 계절 전체 ERA5-Land 기후로 계산하였다. 점은 라벨 셀, 옅은 색은 적용 가능 영역 밖의 셀, 회색 육지는 지도 범위 밖(ESA CCI 영구동토 확률 50\% 미만, 물, 입력 결측)이다. (b) 2011--2024년 각 해의 평균 대비 편차. 모든 해에 같은 색 범위($\pm$15 cm)를 쓴다. (c) 알래스카와 레나 델타의 예측 지역 평균 편차(선, 2011--2024년)와 셀-연도 50개 이상인 해의 관측 셀 평균 편차(원; 알래스카 \nanobsyrsAk{}년, 레나 델타 \nanobsyrsLe{}년). 1 km 값은 표시용으로 250 m에 보간하였다.
